% !TEX root = ../Main.tex
\chapter{补形法：补成长方体}

把三棱锥补成长方体后，外接球就是长方体的外接球：球心在长方体中心，直径就是空间对角线，不必再单独找球心。

\section{对棱相等的三棱锥}

\state{对棱相等 $\to$ 长方体}{
    若三棱锥三组对棱两两相等，记 $AB=CD=a$、$BC=AD=c$、$AC=BD=b$，则可把它嵌进一个长方体，让\textbf{三棱锥的六条棱成为长方体的六条面对角线}。设长方体棱长 $m,n,p$：
    \[
    \begin{cases} m^2+p^2=a^2\\ n^2+p^2=c^2\\ m^2+n^2=b^2 \end{cases}
    \ \Longrightarrow\ 2(m^2+n^2+p^2)=a^2+b^2+c^2.
    \]
    长方体空间对角线就是外接球直径，$m^2+n^2+p^2=(2R)^2$，故
    \[
    R^2=\frac{a^2+b^2+c^2}{8}.
    \]
}

\begin{center}
\begin{tikzpicture}[scale=0.9]
    \draw[thick] (0,0) -- (4,0) -- (4,2) -- (0,2) -- cycle;
    \draw[thick] (1,1) -- (5,1) -- (5,3) -- (1,3) -- cycle;
    \draw[thick] (0,0)--(1,1) (4,0)--(5,1) (4,2)--(5,3) (0,2)--(1,3);
    \coordinate (V1) at (0,0);
    \coordinate (V2) at (5,1);
    \coordinate (V3) at (4,2);
    \coordinate (V4) at (1,3);
    \draw[thick, red] (V1)--(V2) (V1)--(V3) (V1)--(V4) (V2)--(V3) (V2)--(V4) (V3)--(V4);
    \node at (2.5,-0.6) {\small 对棱相等的三棱锥：六条棱就是六条面对角线};
\end{tikzpicture}
\end{center}

\section{墙角三棱锥}

\state{三棱两两垂直 $\to$ 长方体}{
    \textbf{墙角三棱锥}：一个顶点处有三条棱 $a,b,c$ 两两垂直，像墙角一样。把它补回长方体，这三条棱就是长方体过一个顶点的三条棱。空间对角线 $=2R$：
    \[
    a^2+b^2+c^2=(2R)^2\ \Longrightarrow\ R^2=\frac{a^2+b^2+c^2}{4}.
    \]
}

\begin{center}
\begin{tikzpicture}[scale=0.95]
    \coordinate (P) at (0,0);
    \coordinate (A) at (3,0);
    \coordinate (B) at (0,2.4);
    \coordinate (C) at (1.1,1.3);
    \draw[thick] (P)--(A) node[midway,below]{\small $a$};
    \draw[thick] (P)--(B) node[midway,left]{\small $b$};
    \draw[thick] (P)--(C) node[midway,below right]{\small $c$};
    \draw[thick] (A)--(B) (A)--(C) (B)--(C);
    \draw[dashed,gray] (A)--($(A)+(C)$) (B)--($(B)+(C)$) (C)--($(C)+(A)$);
    \draw[dashed,gray] ($(A)+(B)$)--($(A)+(B)+(C)$);
    \draw[dashed,gray] (A)--($(A)+(B)$)--(B);
    \draw[dashed,gray] ($(A)+(C)$)--($(A)+(B)+(C)$)--($(B)+(C)$);
    \draw[thick,red,dashed] (P)--($(A)+(B)+(C)$) node[midway,above]{\small $2R$};
    \node at (P)[below left]{$P$};
    \node at (A)[right]{$A$};
    \node at (B)[above left]{$B$};
    \node at (C)[above]{$C$};
\end{tikzpicture}
\end{center}

\section{墙角模型的"复体"}

\state{以三条面对角线为棱 $\to$ 长方体}{
    有时三棱锥的三条棱 $m,n,p$ 不是长方体的棱，而是它三个面的\textbf{对角线}——把墙角模型"复制"一份拼成的样子。设长方体棱长 $a,b,c$：
    \[
    \begin{cases} m^2=a^2+b^2\\ n^2=b^2+c^2\\ p^2=a^2+c^2 \end{cases}
    \ \Longrightarrow\ m^2+n^2+p^2=2(a^2+b^2+c^2).
    \]
    又 $a^2+b^2+c^2=(2R)^2$，故
    \[
    R^2=\frac{m^2+n^2+p^2}{8}.
    \]
    与对棱相等公式同形（都除以 $8$），区别只在 $m,n,p$ 是哪种棱：对棱相等时是三棱锥的实际棱长，这里是"面对角线型"的棱长。
}
